\documentclass[11pt]{article}
\usepackage{amsmath}
\begin{document}

\TITLE{Rain or Shine: Does Weather Shape Bike-Share Demand?\\Pricing and Fleet Rebalancing}
\KEYWORDS{bike sharing; weather shocks; demand elasticity; fleet rebalancing}

\ABSTRACT{Riders who skip a trip on a rainy day may still shift their later
travel plans. We study how short-run weather shocks affect a platform's
pricing and fleet rebalancing in a two-period model. The operator commits to
a fixed price and a rebalancing capacity. Under the maintained equilibrium,
the model delivers a three-region dispatch rule: static pricing in mild
weather, dynamic rebalancing after either demand shock, and capacity
expansion when weather sensitivity exceeds a threshold.}

\section{Introduction}
Platforms increasingly sell mobility under uncertainty. Dockless bike-share
systems ask riders to commit before learning whether a nearby bike is
available. These systems make demand swings visible: apps display nearby
bikes, and operators publish heat maps of trip origins.

\section{Model}
Two riders observe each other's trip outcomes over two periods. A trip
succeeds with probability $q$ and fails (no bike nearby) otherwise.
The operator commits to a fixed price $p$ and a rebalancing capacity $K$.

\section{Static Weather Response}
When it rained in the morning, the evening rider suffers a scheduling cost
$c$ after a failed trip, which motivates deferred trips. Dynamic rebalancing
is optimal when the demand gap and weather sensitivity are high.

\section{Dynamic Weather Response}
Conditional on a failed trip in period 1, a rider faces state $(0,1)$ if her
neighbor succeeded, or $(0,0)$ if neither did. A fixed price induces one of
four response profiles: UR, CR, NR, or RR.

\section{Managerial Insights}
Stronger weather sensitivity can make rebalancing capacity more widely
valuable and eventually eliminate flat pricing. The operator does not choose
riders' weather sensitivity.

\section{Concluding Remarks}
Persistent exposure can support higher prices without creating additional
trip value. Broader fleet coverage gives riders higher welfare.

\end{document}
